\documentclass{beamer}

\usepackage{beamerthemeSingapore}
\usepackage[brazil]{babel}
\usepackage[T1]{fontenc}
\usepackage[utf8]{inputenc}
\usepackage{graphicx}
\usepackage{url}
\usepackage{xcolor}
\usepackage{booktabs}
\usepackage{ragged2e}
\usepackage{etoolbox}
\usepackage{amsmath}
\usepackage{minted}
\usepackage{tikz}
\usetikzlibrary{positioning, arrows.meta, shapes.geometric}

\newcommand{\tcb}[1]{\textcolor{blue}{#1}}
\newcommand{\tcr}[1]{\textcolor{red}{#1}}
\newcommand{\tcg}[1]{\textcolor{green!50!black}{#1}}

\tikzset{
  treenode/.style = {align=center, inner sep=0pt, text centered,
    font=\sffamily\small},
  arn_b/.style = {treenode, circle, black, draw=black,
    fill=white, text width=1.7em},
  arn_d/.style = {treenode, circle, white, draw=black,
    fill=blue!60, text width=1.7em},
  arn_x/.style = {treenode, rectangle, draw=black, fill=black!10,
    minimum width=0.6em, minimum height=0.6em}
}

\setminted{
  fontsize=\footnotesize,
  breaklines=true,
  breakanywhere=true,
  frame=lines,
  framesep=2mm,
  baselinestretch=1.05,
  tabsize=2
}

\apptocmd{\frame}{}{\justifying}{}

\setbeamertemplate{navigation symbols}{}

\title{05 - Árvores e percursos}
\author{Prof. Alex Kutzke}
\date{Setembro de 2026}

\begin{document}

% ==========================================
\begin{frame}
  \begin{center}
    \large{Setor de Educação Profissional e Tecnológica}\\
    Tecnologia em Análise e Desenvolvimento de Sistemas\\
    \vspace{1cm}
    \small{DS143 - Estruturas de Dados II}
  \end{center}
  \titlepage
\end{frame}

% ==========================================
\begin{frame}{Objetivos da aula}
  Ao final desta aula você deve ser capaz de:
  \begin{enumerate}
    \item Enunciar a definição recursiva de árvore e aplicar a terminologia;
    \item Representar uma árvore binária em Go com uma \texttt{struct} e dois
          ponteiros;
    \item Escrever os percursos em pré-ordem, in-ordem e pós-ordem, e produzir à
          mão a sequência de visitas de cada um;
    \item Escolher o percurso adequado a uma tarefa, justificando pela ordem em
          que a informação fica disponível;
    \item Implementar o percurso em largura com uma fila;
    \item Relacionar altura, número de nós e custo de alcançar um nó.
  \end{enumerate}

  \vspace{0.3cm}
  \footnotesize Este conteúdo cai na Prova 1.
\end{frame}

% ==========================================
\begin{frame}{Roteiro}
  \tableofcontents
\end{frame}

% ==========================================
\section{Motivação e definição}
% ==========================================

\begin{frame}[fragile]{Uma hierarquia comum}
  \begin{minted}{text}
$ tree --charset=ascii 05/demo

05/demo
|-- curso
|   |-- ds122
|   |   `-- aula_01.txt
|   `-- ds143
|       |-- aula_01.txt
|       `-- aula_02.txt
`-- leia-me.txt
  \end{minted}

  \vspace{0.2cm}
  A regra que descreve \texttt{05/demo} é a mesma que descreve
  \texttt{05/demo/curso} e a mesma que descreve
  \texttt{05/demo/curso/ds143}.
\end{frame}

\begin{frame}[fragile]{O tamanho de cada item}
  \begin{minted}{text}
$ du -ab 05/demo

30	05/demo/curso/ds143/aula_01.txt
48	05/demo/curso/ds143/aula_02.txt
78	05/demo/curso/ds143
32	05/demo/curso/ds122/aula_01.txt
32	05/demo/curso/ds122
110	05/demo/curso
46	05/demo/leia-me.txt
156	05/demo
  \end{minted}

  \begin{itemize}
    \item \tcb{Ordem}: o diretório aparece depois do que está dentro dele;
    \item \tcb{Aritmética}: $30 + 48 = 78$, $78 + 32 = 110$, $110 + 46 = 156$.
  \end{itemize}
\end{frame}

\begin{frame}{O tamanho do diretório não existe antes do conteúdo}
  \begin{block}{}
    O \texttt{du} não tinha alternativa. O tamanho de um diretório é a soma dos
    tamanhos do que está dentro, e essa soma só pode ser calculada depois que as
    partes forem conhecidas.
  \end{block}

  \vspace{0.6cm}
  Ao final da aula, esses dois exemplos voltam com nomes específicos: o \texttt{du}
  percorre o sistema de arquivos em \tcb{pós-ordem}, e o \texttt{tree} o
  percorre em \tcb{pré-ordem}.
\end{frame}

\begin{frame}{Você já construiu árvores nesta disciplina}
  \textbf{Quick-Union}, na aula de análise de algoritmos:

  \begin{itemize}
    \item \texttt{id[i]} guardava o \tcb{pai} do elemento \texttt{i};
    \item cada componente conexa era uma árvore, e o conjunto delas uma
          floresta;
    \item \texttt{Find} subia da folha até a raiz.
  \end{itemize}

  \vspace{0.4cm}
  A diferença entre o Quick-Union e o Weighted Quick-Union foi a \tcr{altura}
  dessas árvores: pendurar a menor sob a maior mantém a altura em
  $O(\log n)$; pendurar sem critério permite que chegue a $O(n)$.
\end{frame}

\begin{frame}[fragile]{Você já percorreu árvores nesta disciplina}
  \textbf{Parser de expressões prefixadas}, no material de recursão:

  \begin{minted}{go}
if pos < len(expr) && expr[pos] == '+' {
    pos++
    return eval() + eval()
}
  \end{minted}

  \vspace{0.3cm}
  \begin{itemize}
    \item cada chamada de \texttt{eval} tratava uma subexpressão;
    \item as chamadas aninhadas formavam uma árvore;
    \item ali a árvore \textbf{nunca foi construída na memória}: existia apenas
          como o desenho das chamadas recursivas.
  \end{itemize}
\end{frame}

\begin{frame}{Definição}
  Uma árvore é um conjunto finito de $n$ nós, tal que, se $n > 0$:

  \begin{itemize}
    \item existe um nó distinguido $r$, chamado \tcb{raiz} da árvore;
    \item os $n-1$ nós restantes são divididos em $m \geq 0$ conjuntos
          disjuntos, cada um deles também uma árvore, chamados \tcb{subárvores}
          de $r$.
  \end{itemize}

  \vspace{0.5cm}
  A definição é \tcr{recursiva}, e é assim que vamos programar todas as
  operações de hoje. O caso base é $n = 0$: a \textbf{árvore vazia}, que não tem
  raiz.
\end{frame}

\begin{frame}{Duas consequências da definição}
  \begin{block}{Existe exatamente um caminho da raiz até cada nó}
    Se houvesse dois, algum nó teria dois pais, e os conjuntos deixariam de ser
    disjuntos. A estrutura ainda seria útil, mas seria um \textit{grafo},
    assunto do final do semestre.
  \end{block}

  \begin{block}{A definição não distingue árvore de subárvore}
    Toda subárvore é, também, uma árvore completa. Uma função que
    recebe uma árvore pode ser chamada sobre uma subárvore sem nenhuma
    adaptação.
  \end{block}
\end{frame}

% ==========================================
\section{Terminologia}
% ==========================================

\begin{frame}{A árvore de exemplo}
  \begin{center}
    \begin{tikzpicture}[-, level distance=1.1cm,
        level 1/.style={sibling distance=4cm},
        level 2/.style={sibling distance=2cm},
        level 3/.style={sibling distance=1.4cm}]
      \node [arn_b] {50}
      child { node [arn_b] {30}
          child { node [arn_b] {20}
              child { node [arn_b] {10} }
              child [missing]
            }
          child { node [arn_b] {40}
              child { node [arn_b] {35} }
              child { node [arn_b] {45} }
            }
        }
      child { node [arn_b] {90}
          child [missing]
          child { node [arn_b] {95} }
        };
    \end{tikzpicture}
  \end{center}

  \vspace{0.2cm}
  Desenhamos a raiz em cima e as folhas embaixo. As ligações apontam sempre do
  pai para os filhos, e por isso as setas não são desenhadas.
\end{frame}

\begin{frame}{Raiz, pai, filho, folha, nó interno}
  \begin{columns}
    \begin{column}{0.45\textwidth}
      \centering
      \begin{tikzpicture}[-, level distance=0.9cm,
          level 1/.style={sibling distance=2.6cm},
          level 2/.style={sibling distance=1.3cm},
          level 3/.style={sibling distance=1cm},
          every node/.style={font=\sffamily\scriptsize}]
        \node [arn_b] {50}
        child { node [arn_b] {30}
            child { node [arn_b] {20}
                child { node [arn_d] {10} }
                child [missing]
              }
            child { node [arn_b] {40}
                child { node [arn_d] {35} }
                child { node [arn_d] {45} }
              }
          }
        child { node [arn_b] {90}
            child [missing]
            child { node [arn_d] {95} }
          };
      \end{tikzpicture}
    \end{column}
    \begin{column}{0.55\textwidth}
      \small
      \begin{description}
        \item[Raiz] o único nó sem pai: 50
        \item[Pai de $x$] o nó imediatamente acima: o pai de 35 é 40
        \item[Filho de $x$] nó imediatamente abaixo: 20 e 40 são filhos de 30
        \item[Folha] nó sem filhos: \tcb{10, 35, 45, 95}
        \item[Nó interno] nó com pelo menos um filho: 50, 30, 20, 40, 90
      \end{description}
    \end{column}
  \end{columns}
\end{frame}

\begin{frame}{Grau, nível e altura}
  \begin{description}
    \item[Grau de um nó] número de subárvores do nó. Grau de 30 é 2, de 20 é 1,
          de 10 é 0;
    \item[Grau da árvore] maior grau entre todos os nós. Aqui, 2;
    \item[Nível (profundidade)] número de arestas da raiz até o nó. 50 está no
          nível 0, e 35 no nível 3;
    \item[Altura] comprimento do caminho mais longo da raiz até uma folha. Aqui,
          3;
    \item[Floresta] conjunto de árvores disjuntas, como as componentes do
          Quick-Union.
  \end{description}

  \vspace{0.3cm}
  \begin{block}{Duas convenções desta disciplina}
    A altura é contada em \tcr{arestas}: a árvore com um único nó tem altura 0.
    A árvore \tcr{vazia tem altura $-1$}, valor escolhido para que
    $1 + \max(\text{esq}, \text{dir})$ devolva 0 para uma folha.
  \end{block}
\end{frame}

\begin{frame}{Árvores binárias}
  Árvore de grau máximo 2, com uma exigência a mais: cada subárvore é
  identificada como sendo a da \tcb{esquerda} ou a da \tcb{direita}. Qualquer
  uma das duas pode ser vazia.

  \vspace{0.5cm}
  \begin{columns}
    \begin{column}{0.5\textwidth}
      \centering
      \begin{tikzpicture}[-, level distance=1cm]
        \node [arn_b] {a}
        child { node [arn_b] {b} }
        child [missing];
      \end{tikzpicture}
    \end{column}
    \begin{column}{0.5\textwidth}
      \centering
      \begin{tikzpicture}[-, level distance=1cm]
        \node [arn_b] {a}
        child [missing]
        child { node [arn_b] {b} };
      \end{tikzpicture}
    \end{column}
  \end{columns}

  \vspace{0.5cm}
  Iguais enquanto árvores de grau 2, \tcr{diferentes} enquanto árvores binárias.
  A partir da próxima aula a posição do filho carrega informação: à esquerda o
  que é menor, à direita o que é maior.
\end{frame}

% ==========================================
\section{Representação em Go}
% ==========================================

\begin{frame}[fragile]{A estrutura de um nó}
  \begin{minted}{go}
type No struct {
    Info int
    Esq  *No
    Dir  *No
}
  \end{minted}

  \vspace{0.3cm}
  \begin{itemize}
    \item O tipo se refere a si mesmo, o que só é possível através de
          \tcb{ponteiros}. Um campo \texttt{Esq No} daria um tipo de tamanho
          infinito, e o compilador recusa
          (\texttt{invalid recursive type No});
    \item A árvore inteira é o \tcb{ponteiro para o nó raiz}. Um \texttt{*No} é
          ao mesmo tempo um nó e a subárvore que começa nele;
    \item A árvore vazia é \texttt{nil}.
  \end{itemize}
\end{frame}

\begin{frame}[fragile]{Sem criar vazia, sem liberar}
  Em C seriam necessárias duas funções a mais:

  \begin{minted}{c}
Arvore* cria_arv_vazia (void) { return NULL; }

Arvore* arv_libera (Arvore* a) {
  if (a != NULL) {
    arv_libera(a->esq);
    arv_libera(a->dir);
    free(a);
  }
  return NULL;
}
  \end{minted}

  \vspace{0.2cm}
  Em Go, a árvore vazia é \texttt{nil} e o coletor de lixo recolhe os nós que
  ficam sem referência. \footnotesize Note, de passagem, a ordem das chamadas na
  \texttt{arv\_libera}: os filhos antes do \texttt{free} do pai.
\end{frame}

\begin{frame}[fragile]{Duas funções de construção}
  \begin{minted}{go}
func Constroi(info int, esq, dir *No) *No {
    return &No{Info: info, Esq: esq, Dir: dir}
}

func Folha(info int) *No {
    return Constroi(info, nil, nil)
}
  \end{minted}

  \vspace{0.3cm}
  \texttt{\&No\{...\}} cria o nó e devolve o endereço dele. Em Go isso é seguro
  mesmo com a variável saindo de escopo: o compilador percebe que o endereço
  escapa da função e aloca o valor onde ele sobreviva.
\end{frame}

\begin{frame}[fragile]{A árvore de exemplo em uma expressão}
  \begin{minted}{go}
func ArvoreExemplo() *No {
    return Constroi(50,
        Constroi(30,
            Constroi(20,
                Folha(10),
                nil),
            Constroi(40,
                Folha(35),
                Folha(45))),
        Constroi(90,
            nil,
            Folha(95)))
}
  \end{minted}

  \vspace{0.2cm}
  O recuo do código ajuda a enxergar o desenho da árvore.
\end{frame}

\begin{frame}[fragile]{Procurar um valor}
  \begin{minted}{go}
func Pertence(a *No, procurado int) bool {
    if a == nil {
        return false
    }
    if a.Info == procurado {
        return true
    }
    return Pertence(a.Esq, procurado) ||
        Pertence(a.Dir, procurado)
}
  \end{minted}

  \vspace{0.3cm}
  No pior caso, visita todos os $n$ nós, e não há como fazer melhor: nada nesta
  árvore diz de que lado procurar. Essa é a lacuna que a \tcb{árvore binária de
    busca} preenche na próxima aula.
\end{frame}

% ==========================================
\section{Percursos em profundidade}
% ==========================================

\begin{frame}{Percorrer e visitar}
  \begin{block}{}
    \textbf{Percorrer} uma árvore é visitar sistematicamente cada um dos seus
    nós, uma vez cada. \textbf{Visitar} um nó é fazer alguma coisa com a
    informação dele: imprimir, somar, alterar, comparar.
  \end{block}

  \vspace{0.4cm}
  Durante um percurso é comum \textit{passar} por um nó várias vezes sem
  \textit{visitá-lo}.

  \vspace{0.4cm}
  Para aplicar um reajuste de 5\% em todos os salários guardados na árvore,
  qualquer ordem serve. Para o \texttt{du} da abertura, a ordem era o problema
  inteiro.
\end{frame}

\begin{frame}{Três ordens, uma diferença}
  Os três percursos em profundidade descem por um ramo até o fim antes de tentar
  o ramo seguinte. Diferem apenas em \tcr{quando a raiz é visitada}:

  \vspace{0.5cm}
  \begin{center}
    \begin{tabular}{lll}
      \toprule
      \textbf{Percurso} & \textbf{Ordem}           & \textbf{Sigla} \\
      \midrule
      Pré-ordem         & raiz, esquerda, direita  & R, E, D        \\
      In-ordem          & esquerda, raiz, direita  & E, R, D        \\
      Pós-ordem         & esquerda, direita, raiz  & E, D, R        \\
      \bottomrule
    \end{tabular}
  \end{center}

  \vspace{0.5cm}
  Os três códigos têm as mesmas três linhas, em ordens diferentes.
\end{frame}

\begin{frame}[fragile]{Pré-ordem (R, E, D)}
  \begin{columns}
    \begin{column}{0.55\textwidth}
      \begin{minted}{go}
func PreOrdem(a *No) {
    if a == nil {
        return
    }
    fmt.Print(a.Info, " ")
    PreOrdem(a.Esq)
    PreOrdem(a.Dir)
}
      \end{minted}
    \end{column}
    \begin{column}{0.45\textwidth}
      \centering
      \begin{tikzpicture}[-, level distance=0.85cm,
          level 1/.style={sibling distance=2.2cm},
          level 2/.style={sibling distance=1.1cm},
          level 3/.style={sibling distance=0.9cm},
          every node/.style={font=\sffamily\scriptsize}]
        \node [arn_b] {50}
        child { node [arn_b] {30}
            child { node [arn_b] {20}
                child { node [arn_b] {10} }
                child [missing]
              }
            child { node [arn_b] {40}
                child { node [arn_b] {35} }
                child { node [arn_b] {45} }
              }
          }
        child { node [arn_b] {90}
            child [missing]
            child { node [arn_b] {95} }
          };
      \end{tikzpicture}
    \end{column}
  \end{columns}

  \vspace{0.3cm}
  \texttt{50 30 20 10 40 35 45 90 95}

  \footnotesize O \texttt{if a == nil} é o caso base, e corresponde à árvore
  vazia da definição. Sem ele o programa quebra na primeira folha.
\end{frame}

\begin{frame}[fragile]{In-ordem (E, R, D)}
  \begin{columns}
    \begin{column}{0.55\textwidth}
      \begin{minted}{go}
func InOrdem(a *No) {
    if a == nil {
        return
    }
    InOrdem(a.Esq)
    fmt.Print(a.Info, " ")
    InOrdem(a.Dir)
}
      \end{minted}
    \end{column}
    \begin{column}{0.45\textwidth}
      \centering
      \begin{tikzpicture}[-, level distance=0.85cm,
          level 1/.style={sibling distance=2.2cm},
          level 2/.style={sibling distance=1.1cm},
          level 3/.style={sibling distance=0.9cm},
          every node/.style={font=\sffamily\scriptsize}]
        \node [arn_b] {50}
        child { node [arn_b] {30}
            child { node [arn_b] {20}
                child { node [arn_b] {10} }
                child [missing]
              }
            child { node [arn_b] {40}
                child { node [arn_b] {35} }
                child { node [arn_b] {45} }
              }
          }
        child { node [arn_b] {90}
            child [missing]
            child { node [arn_b] {95} }
          };
      \end{tikzpicture}
    \end{column}
  \end{columns}

  \vspace{0.3cm}
  \texttt{10 20 30 35 40 45 50 90 95}

  \footnotesize A saída veio em ordem crescente. Isso é propriedade
  \textbf{desta} árvore, montada com os menores à esquerda, e não de toda árvore
  binária. A próxima aula parte desse fato.
\end{frame}

\begin{frame}[fragile]{Pós-ordem (E, D, R)}
  \begin{columns}
    \begin{column}{0.55\textwidth}
      \begin{minted}{go}
func PosOrdem(a *No) {
    if a == nil {
        return
    }
    PosOrdem(a.Esq)
    PosOrdem(a.Dir)
    fmt.Print(a.Info, " ")
}
      \end{minted}
    \end{column}
    \begin{column}{0.45\textwidth}
      \centering
      \begin{tikzpicture}[-, level distance=0.85cm,
          level 1/.style={sibling distance=2.2cm},
          level 2/.style={sibling distance=1.1cm},
          level 3/.style={sibling distance=0.9cm},
          every node/.style={font=\sffamily\scriptsize}]
        \node [arn_b] {50}
        child { node [arn_b] {30}
            child { node [arn_b] {20}
                child { node [arn_b] {10} }
                child [missing]
              }
            child { node [arn_b] {40}
                child { node [arn_b] {35} }
                child { node [arn_b] {45} }
              }
          }
        child { node [arn_b] {90}
            child [missing]
            child { node [arn_b] {95} }
          };
      \end{tikzpicture}
    \end{column}
  \end{columns}

  \vspace{0.3cm}
  \texttt{10 20 35 45 40 30 95 90 50}

  \footnotesize É o percurso do \texttt{du}: 10 e 20 antes de 30, e a raiz 50 no
  fim de tudo.
\end{frame}

\begin{frame}{Escolhendo o percurso pela tarefa}
  \small
  \begin{center}
    \begin{tabular}{p{3.6cm}p{2cm}p{4.4cm}}
      \toprule
      \textbf{Tarefa}                                  & \textbf{Percurso} & \textbf{Por quê}                                                  \\
      \midrule
      Copiar a árvore preservando a forma              & pré-ordem         & a raiz precisa existir antes que os filhos sejam pendurados nela  \\
      Listar os valores de uma ABB em ordem            & in-ordem          & é a ordem que a organização da ABB garante                        \\
      Somar tamanhos, calcular altura, liberar memória & pós-ordem         & o resultado do nó depende dos resultados dos filhos               \\
      Achar o nó mais próximo da raiz                  & em largura        & os níveis são visitados em ordem crescente de profundidade        \\
      \bottomrule
    \end{tabular}
  \end{center}

  \vspace{0.3cm}
  A pós-ordem é o caso mais rígido: trocar a ordem das linhas
  \tcr{quebra a função}.
\end{frame}

\begin{frame}{Árvore de expressão}
  O parser do material de recursão, agora com a árvore construída de fato:
  $2 \times (3 + 4)$.

  \vspace{0.4cm}
  \begin{center}
    \begin{tikzpicture}[-, level distance=1.1cm,
        level 1/.style={sibling distance=2.6cm},
        level 2/.style={sibling distance=1.6cm}]
      \node [arn_b] {*}
      child { node [arn_b] {2} }
      child { node [arn_b] {+}
          child { node [arn_b] {3} }
          child { node [arn_b] {4} }
        };
    \end{tikzpicture}
  \end{center}

  \vspace{0.3cm}
  Operadores nos nós internos, números nas folhas.
\end{frame}

\begin{frame}[fragile]{Os três percursos dão as três notações}
  \begin{minted}{text}
pré-ordem (prefixa)       : * 2 + 3 4
in-ordem  (sem parênteses): 2 * 3 + 4
in-ordem  (com parênteses): (2 * (3 + 4))
pós-ordem (pós-fixa)      : 2 3 4 + *
valor                     : 14
  \end{minted}

  \vspace{0.3cm}
  \begin{itemize}
    \item A primeira linha é exatamente a entrada que o
          \texttt{prefix\_parser.go} lê;
    \item Prefixa e pós-fixa dispensam parênteses: a posição do operador já
          determina a estrutura;
    \item O in-ordem puro perde essa informação. \texttt{2 * 3 + 4} lido com a
          precedência usual vale \tcr{10}, e não 14.
  \end{itemize}
\end{frame}

\begin{frame}[fragile]{Avaliar é um percurso em pós-ordem}
  \begin{minted}{go}
func Avalia(a *No) int {
    if a.Esq == nil && a.Dir == nil {
        valor, _ := strconv.Atoi(a.Info)
        return valor
    }

    esq := Avalia(a.Esq)
    dir := Avalia(a.Dir)

    switch a.Info {
    case "+":
        return esq + dir
    case "*":
        return esq * dir
    }
    return 0
}
  \end{minted}
\end{frame}

% ==========================================
\section{Percurso em largura}
% ==========================================

\begin{frame}[fragile]{Nível a nível}
  Os três percursos anteriores descem até o fim de um ramo antes de olhar o
  seguinte. O percurso em largura visita todos os nós do nível 0, depois todos
  os do nível 1, e assim por diante.

  \vspace{0.4cm}
  \begin{minted}{text}
em largura: 50 30 90 20 40 95 10 35 45
pré-ordem : 50 30 20 10 40 35 45 90 95
  \end{minted}

  \vspace{0.3cm}
  A pré-ordem alcança o 10, que está no nível 3, antes de alcançar o 90, que
  está no nível 1.
\end{frame}

\begin{frame}{Por que a recursão não resolve}
  \begin{block}{}
    A recursão dá acesso direto aos \textbf{filhos} de quem está sendo
    visitado. O percurso em largura precisa dos \textbf{irmãos}, que estão em
    outro ramo.
  \end{block}

  \vspace{0.5cm}
  A solução é guardar explicitamente os nós ainda por visitar, em uma
  \tcb{fila}: retira-se um nó da frente, visita-se, e seus filhos entram no fim.
\end{frame}

\begin{frame}[fragile]{Percurso em largura com uma fila}
  \begin{minted}[fontsize=\scriptsize]{go}
func EmLargura(a *No) {
    if a == nil {
        return
    }

    fila := []*No{a}

    for len(fila) > 0 {
        atual := fila[0]
        fila = fila[1:]

        fmt.Print(atual.Info, " ")

        if atual.Esq != nil {
            fila = append(fila, atual.Esq)
        }
        if atual.Dir != nil {
            fila = append(fila, atual.Dir)
        }
    }
}
  \end{minted}
\end{frame}

\begin{frame}[fragile]{Trocar a fila por uma pilha devolve a pré-ordem}
  \begin{minted}[fontsize=\scriptsize]{go}
pilha := []*No{a}

for len(pilha) > 0 {
    topo := pilha[len(pilha)-1]
    pilha = pilha[:len(pilha)-1]

    fmt.Print(topo.Info, " ")

    if topo.Dir != nil {
        pilha = append(pilha, topo.Dir)
    }
    if topo.Esq != nil {
        pilha = append(pilha, topo.Esq)
    }
}
  \end{minted}

  \footnotesize O filho da direita é empilhado primeiro para que o da esquerda
  saia antes. É o que a recursão vinha fazendo na pilha de chamadas do programa.
\end{frame}

\begin{frame}{Custo dos percursos}
  \begin{itemize}
    \item \textbf{Tempo}: $\Theta(n)$ nos quatro. Cada nó é visitado uma vez e
          cada ponteiro é examinado uma vez;
    \item \textbf{Espaço, em profundidade}: a pilha de chamadas, proporcional à
          \tcb{altura}. $O(\log n)$ na árvore cheia, $O(n)$ na degenerada;
    \item \textbf{Espaço, em largura}: a fila chega a conter um nível inteiro,
          até $n/2$ nós no último nível de uma árvore cheia.
  \end{itemize}

  \vspace{0.5cm}
  O percurso em largura é o mesmo \textbf{BFS} que aparecerá em grafos, no final
  do semestre, com a mesma fila e a mesma estrutura de laço.
\end{frame}

% ==========================================
\section{Altura e desempenho}
% ==========================================

\begin{frame}[fragile]{Calcular a altura}
  \begin{minted}{go}
func Altura(a *No) int {
    if a == nil {
        return -1
    }

    esq := Altura(a.Esq)
    dir := Altura(a.Dir)

    if esq > dir {
        return 1 + esq
    }
    return 1 + dir
}
  \end{minted}

  \vspace{0.2cm}
  O $-1$ da árvore vazia é o que faz uma folha devolver $1 + (-1) = 0$, sem caso
  especial.
\end{frame}

\begin{frame}{Árvore cheia}
  Todo nó interno tem duas subárvores e todas as folhas estão no último nível.

  \vspace{0.3cm}
  \begin{center}
    \begin{tikzpicture}[-, level distance=0.9cm,
        level 1/.style={sibling distance=3cm},
        level 2/.style={sibling distance=1.5cm}]
      \node [arn_b] {}
      child { node [arn_b] {}
          child { node [arn_b] {} }
          child { node [arn_b] {} }
        }
      child { node [arn_b] {}
          child { node [arn_b] {} }
          child { node [arn_b] {} }
        };
    \end{tikzpicture}
  \end{center}

  \vspace{0.3cm}
  No nível $k$ há $2^k$ nós. Somando todos os níveis de uma árvore cheia de
  altura $h$:

  \[ n = 2^0 + 2^1 + \cdots + 2^h = 2^{h+1} - 1 \]

  Isolando $h$: uma árvore cheia com $n$ nós tem altura
  $h = \log_2(n+1) - 1$, ou seja, $\Theta(\log n)$.
\end{frame}

\begin{frame}{Árvore degenerada}
  Um único filho por nó interno, um nó por nível. Com $n$ nós, a altura é
  $n - 1$.

  \vspace{0.3cm}
  \begin{center}
    \begin{tikzpicture}[-, level distance=0.8cm]
      \node [arn_b] {}
      child { node [arn_b] {}
          child { node [arn_b] {}
              child { node [arn_b] {} }
              child [missing]
            }
          child [missing]
        }
      child [missing];
    \end{tikzpicture}
  \end{center}

  \vspace{0.3cm}
  A forma é a de uma \tcr{lista encadeada}, e o desempenho também.
\end{frame}

\begin{frame}[fragile]{Mesmo número de nós, duas formas}
  \begin{minted}{text}
$ go run propriedades.go

     nós |    cheia | degenerada
---------+----------+-----------
      15 |        3 |         14
      63 |        5 |         62
     255 |        7 |        254
    1023 |        9 |       1022
  \end{minted}

  \vspace{0.3cm}
  Com 1023 nós, o caminho até o nó mais distante é mais de cem vezes mais longo
  na árvore degenerada, com a mesma quantidade de informação armazenada.
\end{frame}

\begin{frame}{Altura mínima e altura máxima}
  \begin{block}{}
    Uma árvore binária com $n$ nós tem altura mínima proporcional a $\log_2 n$,
    no caso da árvore cheia, e altura máxima proporcional a $n$, no caso da
    degenerada. A altura indica o esforço para alcançar qualquer nó.
  \end{block}

  \vspace{0.5cm}
  É a mesma diferença entre o Quick-Union e o Weighted Quick-Union, e pela mesma
  razão. Garantir que a árvore fique perto do primeiro caso, e não do segundo, é
  o assunto das duas próximas aulas.
\end{frame}

\begin{frame}{De volta ao \texttt{tree} e ao \texttt{du}}
  \begin{itemize}
    \item O \texttt{tree} imprime cada diretório \textbf{antes} do seu conteúdo:
          percorre em \tcb{pré-ordem}. Quem lê precisa saber em que diretório
          está antes de ver os arquivos;
    \item O \texttt{du} imprime cada diretório \textbf{depois} do seu conteúdo:
          percorre em \tcb{pós-ordem}. O número do diretório é a soma dos
          números dos filhos, e não existe antes deles.
  \end{itemize}

  \vspace{0.5cm}
  O sistema de arquivos não é uma árvore binária: um diretório tem quantos
  filhos quiser. É o caso geral da definição, com $m$ subárvores. Os percursos
  continuam valendo, com um laço sobre os filhos no lugar das duas chamadas
  fixas.
\end{frame}

\begin{frame}{Exercícios}
  Os exercícios estão no texto da aula. \textbf{Não valem nota} e não têm
  entrega: são a preparação para a aula prática de árvores, essa sim avaliada.
  O conteúdo cai na Prova 1.

  \vspace{0.4cm}
  \begin{itemize}
    \item grau, folhas, altura e nível sobre uma árvore dada;
    \item as quatro sequências de percurso, à mão;
    \item \texttt{ContaNos}, \texttt{Maximo}, \texttt{Espelha},
          \texttt{Iguais}, \texttt{NoNoNivel};
    \item impressão na notação \texttt{<raiz esq dir>};
    \item desafio: reconstruir a árvore a partir da pré-ordem e do in-ordem.
  \end{itemize}

  \vspace{0.4cm}
  Resolva antes da próxima aula, que começa pela árvore binária de busca.
\end{frame}

\end{document}
